% !TEX program = lualatex
% 機械加工・組立・溶接部品 要求仕様書テンプレート
% 推奨コンパイル: lualatex spec.tex

\documentclass[a4paper,10.5pt]{ltjsarticle}

% ===== Page / layout =====
\usepackage[top=18mm,bottom=18mm,left=16mm,right=16mm]{geometry}
\usepackage{array,tabularx,longtable,multirow}
\usepackage{enumitem}
\usepackage[table]{xcolor}
\usepackage[explicit]{titlesec}
\usepackage{fancyhdr}
\usepackage{lastpage}
\usepackage{booktabs}
\usepackage{pifont}
\usepackage{setspace}
\usepackage{needspace}
% 図面番号・仕様書番号などの半角アンダースコア _ を本文中で安全に扱う。

% ===== Color palette: same hue, low saturation =====
\definecolor{SpecInk}{HTML}{1F2F3D}
\definecolor{SpecHead}{HTML}{2D4B63}
\definecolor{SpecSoft}{HTML}{EEF3F6}
\definecolor{SpecLine}{HTML}{CFD8DF}
\definecolor{SpecMuted}{HTML}{5F7180}

\pagecolor{white}
\color{SpecInk}
\setstretch{1.12}
\setlength{\parindent}{0pt}
\setlength{\parskip}{4pt}
\renewcommand{\arraystretch}{1.32}
\setlist[itemize]{leftmargin=1.2em,itemsep=2pt,topsep=3pt}
\setlist[enumerate]{leftmargin=1.6em,itemsep=2pt,topsep=3pt}

% ===== Header / Footer =====
\pagestyle{fancy}
\fancyhf{}
\lhead{\small 機械加工・組立・溶接部品 要求仕様書}
\rhead{\small 仕様書番号：\SpecNo\quad Rev.：\SpecRev}
\lfoot{\small\textcolor{SpecMuted}{\SpecDocTitle}}
\rfoot{\small\textcolor{SpecMuted}{Page \thepage\ / \pageref{LastPage}}}
\renewcommand{\headrulewidth}{0.6pt}
\renewcommand{\footrulewidth}{0.4pt}
\renewcommand{\headrule}{\hbox to\headwidth{\color{SpecHead}\leaders\hrule height \headrulewidth\hfill}}
\renewcommand{\footrule}{\hbox to\headwidth{\color{SpecLine}\leaders\hrule height \footrulewidth\hfill}}

% ===== Title styles =====
\titleformat{\section}
  {\Large\bfseries\color{white}}
  {}{0pt}
  {\colorbox{SpecHead}{\parbox{\dimexpr\linewidth-2\fboxsep}{\thesection\quad #1}}}
\titlespacing*{\section}{0pt}{0pt}{14pt}

\titleformat{\subsection}
  {\large\bfseries\color{SpecHead}}
  {}{0pt}
  {\colorbox{SpecSoft}{\parbox{\dimexpr\linewidth-2\fboxsep}{\textcolor{SpecHead}{\thesubsection\quad #1}}}}
\titlespacing*{\subsection}{0pt}{16pt}{6pt}

% ===== Table helpers =====
\newcolumntype{Y}{>{\raggedright\arraybackslash}X}
\newcolumntype{L}[1]{>{\raggedright\arraybackslash}p{#1}}
\newcommand{\thcell}[1]{\textbf{\textcolor{SpecHead}{#1}}}
\newcommand{\labelcell}[1]{\cellcolor{SpecSoft}\textbf{\textcolor{SpecHead}{#1}}}
\newcommand{\blank}{\rule{0pt}{2.1ex}\hspace{1em}}
\newcommand{\cb}{\ding{113}}
\newcommand{\choice}[1]{\textcolor{SpecMuted}{#1}}

% SpecTable は未使用。tabularx は newenvironment で包まない。

\newenvironment{SpecNote}{%
  \par\vspace{3pt}%
  \begin{center}%
  \begin{minipage}{0.96\linewidth}%
  \small\color{SpecMuted}%
}{%
  \end{minipage}%
  \end{center}%
  \vspace{2pt}%
}

\newcommand{\PageBreak}{\clearpage}

% ===== Editable document variables =====
\newcommand{\SpecDocTitle}{機械加工・組立・溶接部品 要求仕様書}
\newcommand{\SpecNo}{（記入）}
\newcommand{\SpecRev}{（記入）}
\newcommand{\SpecDate}{（記入）年（記入）月（記入）日}
\newcommand{\ProjectName}{（記入）}
\newcommand{\PartName}{（記入）}
\newcommand{\DrawingNo}{（記入）}
\newcommand{\Quantity}{（記入）}
\newcommand{\DueDate}{（記入）年（記入）月（記入）日}

\begin{document}

% 表紙は別途作成。本テンプレートは本文ページ用。

% =========================================================
\section{文書情報・調達対象}

\subsection{文書基本情報}
\begin{tabularx}{\linewidth}{|L{29mm}|Y|L{34mm}|Y|}
\hline
\labelcell{対象品名} & \PartName & \labelcell{適用装置 / 案件名} & \ProjectName \\ \hline
\labelcell{図面番号} & \DrawingNo & \labelcell{図面 Rev.} & （記入） \\ \hline
\labelcell{仕様書番号} & \SpecNo & \labelcell{仕様書 Rev.} & \SpecRev \\ \hline
\labelcell{希望数量} & \Quantity & \labelcell{希望納期} & \DueDate \\ \hline
\labelcell{作成部門 / 作成者} & （記入） & \labelcell{承認者} & （記入） \\ \hline
\labelcell{作成日} & \SpecDate & \labelcell{取引先 / 製作者} & （記入） \\ \hline
\end{tabularx}

\subsection{改訂履歴}
\begin{tabularx}{\linewidth}{|L{15mm}|L{28mm}|Y|L{35mm}|}
\hline
\rowcolor{SpecSoft}\thcell{Rev.} & \thcell{日付} & \thcell{改訂内容} & \thcell{作成 / 承認} \\ \hline
0 & （記入）/（記入）/（記入） & 初版 & （記入） / （記入） \\ \hline
  &  &  &  \\ \hline
  &  &  &  \\ \hline
\end{tabularx}

\subsection{調達対象品}
\begin{tabularx}{\linewidth}{|L{10mm}|Y|L{32mm}|L{18mm}|Y|}
\hline
\rowcolor{SpecSoft}\thcell{No.} & \thcell{品名} & \thcell{図面番号} & \thcell{数量} & \thcell{備考} \\ \hline
1 &  &  &  &  \\ \hline
2 &  &  &  &  \\ \hline
3 &  &  &  &  \\ \hline
4 &  &  &  &  \\ \hline
\end{tabularx}

\PageBreak

% =========================================================
\section{適用範囲・支給資料}

\subsection{本仕様書の適用範囲}
\begin{tabularx}{\linewidth}{|L{38mm}|L{24mm}|Y|}
\hline
\rowcolor{SpecSoft}\thcell{項目} & \thcell{要否} & \thcell{指定内容 / 備考} \\ \hline
材料手配 & \choice{要 / 不要} &  \\ \hline
機械加工 & \choice{要 / 不要} &  \\ \hline
溶接 & \choice{要 / 不要} &  \\ \hline
組立 & \choice{要 / 不要} &  \\ \hline
熱処理 & \choice{要 / 不要} &  \\ \hline
表面処理 & \choice{要 / 不要} &  \\ \hline
検査・検査成績書 & \choice{要 / 不要} &  \\ \hline
梱包・識別・納入 & \choice{要 / 不要} &  \\ \hline
\end{tabularx}

\subsection{対象外範囲}
\begin{tabularx}{\linewidth}{|L{35mm}|Y|}
\hline
\labelcell{対象外作業} &  \\ \hline
\labelcell{別途手配品} &  \\ \hline
\labelcell{支給品} &  \\ \hline
\end{tabularx}

\subsection{図面・仕様の優先順位}
\begin{enumerate}
  \item 本要求仕様書
  \item 承認済み製作図
  \item 部品表
  \item 3D CADデータ
  \item その他参考資料
\end{enumerate}

\begin{SpecNote}
不一致、疑義、製作上の制約を発見した場合は、製作前に発注者へ確認すること。
\end{SpecNote}

\subsection{支給資料}
\begin{tabularx}{\linewidth}{|Y|L{18mm}|L{35mm}|Y|}
\hline
\rowcolor{SpecSoft}\thcell{資料名} & \thcell{Rev.} & \thcell{形式} & \thcell{備考} \\ \hline
製作図 &  & PDF / DXF &  \\ \hline
3D CAD &  & STEP / Parasolid &  \\ \hline
部品表 &  & Excel / PDF &  \\ \hline
検査基準書 &  & PDF & 必要時 \\ \hline
\end{tabularx}

\PageBreak

% =========================================================
\section{材料仕様・加工仕様}

\subsection{材料仕様}
\begin{tabularx}{\linewidth}{|L{30mm}|Y|L{30mm}|Y|}
\hline
\labelcell{材質} & 図面指示による / 下記指定 & \labelcell{材料規格} & JIS / ASTM / EN / その他 \\ \hline
\labelcell{ミルシート} & \choice{要 / 不要} & \labelcell{代替材} & \choice{不可 / 可 / 要承認} \\ \hline
\labelcell{支給材} & \choice{有 / 無} & \labelcell{材料ロット管理} & \choice{要 / 不要} \\ \hline
\end{tabularx}

\begin{SpecNote}
代替材を提案する場合は、機械的性質、耐食性、熱処理性、溶接性、寸法安定性への影響を明記し、発注者の承認を得ること。
\end{SpecNote}

\subsection{部品別材料指定}
\begin{tabularx}{\linewidth}{|L{10mm}|Y|L{30mm}|L{30mm}|Y|}
\hline
\rowcolor{SpecSoft}\thcell{No.} & \thcell{部品名} & \thcell{材質} & \thcell{規格} & \thcell{備考} \\ \hline
1 &  &  &  &  \\ \hline
2 &  &  &  &  \\ \hline
3 &  &  &  &  \\ \hline
\end{tabularx}

\subsection{機械加工の基本要求}
\begin{tabularx}{\linewidth}{|Y|Y|}
\hline
\cb\ 図面指示寸法、公差、幾何公差を満足 & \cb\ 図面指示面粗さを満足 \\ \hline
\cb\ バリ、カエリ、鋭利エッジを除去 & \cb\ ねじ部、はめあい部、シール面を保護 \\ \hline
\cb\ 加工基準変更時は事前確認 & \cb\ 歪み懸念時は加工順序を事前確認 \\ \hline
\end{tabularx}

\subsection{一般公差・標準仕上げ}
\begin{tabularx}{\linewidth}{|L{42mm}|Y|}
\hline
\labelcell{図面指示なき寸法公差} & 例：JIS B 0405-m 相当 / 個別指定 \\ \hline
\labelcell{図面指示なき幾何公差} & 例：JIS B 0419 相当 / 個別指定 \\ \hline
\labelcell{指示なき角部} & 糸面取り / C0.2〜0.5 程度 / 個別指定 \\ \hline
\labelcell{指示なき面粗さ} & 例：Ra 6.3 / 個別指定 \\ \hline
\end{tabularx}

\PageBreak

% =========================================================
\section{溶接仕様・組立仕様}

\subsection{溶接仕様}
\begin{tabularx}{\linewidth}{|L{30mm}|Y|L{30mm}|Y|}
\hline
\labelcell{溶接方法} & TIG / MIG / MAG / アーク / その他 & \labelcell{溶接材料} & 図面指示 / 下記指定 \\ \hline
\labelcell{外観要求} & 割れ、ブローホール、有害なアンダーカットなきこと & \labelcell{歪み管理} & \choice{要 / 不要} \\ \hline
\labelcell{溶接後処理} & スパッタ除去 / ビード仕上げ / 酸洗い & \labelcell{非破壊検査} & 不要 / PT / MT / UT / RT \\ \hline
\end{tabularx}

\subsection{溶接に関する要求事項}
\begin{itemize}
  \item 溶接寸法、脚長、溶接範囲は図面指示によること。
  \item 溶接歪みにより、機能寸法、取付寸法、平面度、直角度を損なわないこと。
  \item 溶接後に機械加工を行う箇所は、必要な取り代を確保すること。
  \item 欠陥補修を行う場合は、補修方法を明確にし、必要に応じて承認を得ること。
\end{itemize}

\subsection{組立仕様}
\begin{tabularx}{\linewidth}{|L{30mm}|Y|L{30mm}|Y|}
\hline
\labelcell{組立範囲} & 全組立 / 部分組立 / 仮組 & \labelcell{締結部品} & 支給 / 製作者手配 \\ \hline
\labelcell{接着剤・固定剤} & 使用可 / 使用不可 / 指定品あり & \labelcell{分解納入} & 可 / 不可 \\ \hline
\labelcell{可動部確認} & 要 / 不要 & \labelcell{締付トルク管理} & 要 / 不要 / 指定部のみ \\ \hline
\end{tabularx}

\subsection{組立に関する要求事項}
\begin{tabularx}{\linewidth}{|Y|Y|}
\hline
\cb\ 部品相互の干渉なし & \cb\ 芯ずれ、傾き、ガタつきなし \\ \hline
\cb\ 可動部の引っ掛かり、異音なし & \cb\ 締結部の緩み止めを確認 \\ \hline
\cb\ 分解困難な構造は事前確認 & \cb\ 組立後の外観・機能を確認 \\ \hline
\end{tabularx}

\PageBreak

% =========================================================
\section{処理仕様・外観仕上げ}

\subsection{熱処理・表面処理}
\begin{tabularx}{\linewidth}{|L{30mm}|Y|L{30mm}|Y|}
\hline
\labelcell{熱処理} & なし / 焼入れ / 焼戻し / 焼鈍 / 応力除去 & \labelcell{処理証明書} & 要 / 不要 \\ \hline
\labelcell{表面処理} & なし / 黒染め / 無電解Ni / アルマイト / 塗装 / その他 & \labelcell{マスキング指定} & あり / なし \\ \hline
\labelcell{硬度指定} & あり / なし & \labelcell{処理後寸法管理} & 要 / 不要 / 重要部のみ \\ \hline
\end{tabularx}

\subsection{処理時の注意事項}
\begin{itemize}
  \item 処理により寸法、公差、ねじ、はめあい、面粗さを損なわないこと。
  \item 処理膜厚または熱変形が重要寸法に影響する場合は、製作前に確認すること。
  \item マスキング箇所、処理不可箇所は図面または別紙で明示すること。
  \item 処理後の錆、変色、剥離、膨れ、密着不良がないこと。
\end{itemize}

\subsection{外観・仕上げ要求}
\begin{tabularx}{\linewidth}{|L{32mm}|Y|}
\hline
\rowcolor{SpecSoft}\thcell{項目} & \thcell{判定基準} \\ \hline
バリ・カエリ & 使用上有害なものがないこと。 \\ \hline
傷・打痕 & 機能面、シール面、摺動面に有害な傷・打痕がないこと。 \\ \hline
錆・変色 & 有害な錆、腐食、処理不良がないこと。 \\ \hline
汚れ・異物 & 切粉、粉塵、溶接スパッタ、過剰な油分がないこと。 \\ \hline
エッジ & 機能上必要な箇所を除き、鋭利でないこと。 \\ \hline
\end{tabularx}

\subsection{使用環境・特記事項}
\begin{tabularx}{\linewidth}{|L{38mm}|Y|}
\hline
\labelcell{使用環境} &  \\ \hline
\labelcell{温度 / 湿度 / 腐食環境} &  \\ \hline
\labelcell{清浄度要求} &  \\ \hline
\end{tabularx}

\PageBreak

% =========================================================
\section{検査要求}

\subsection{検査項目}
\begin{tabularx}{\linewidth}{|Y|L{22mm}|L{24mm}|Y|}
\hline
\rowcolor{SpecSoft}\thcell{検査項目} & \thcell{要否} & \thcell{記録提出} & \thcell{備考} \\ \hline
外観検査 & 要 / 不要 & 要 / 不要 &  \\ \hline
主要寸法検査 & 要 / 不要 & 要 / 不要 &  \\ \hline
幾何公差検査 & 要 / 不要 & 要 / 不要 &  \\ \hline
ねじ検査 & 要 / 不要 & 要 / 不要 &  \\ \hline
表面粗さ検査 & 要 / 不要 & 要 / 不要 &  \\ \hline
硬度検査 & 要 / 不要 & 要 / 不要 &  \\ \hline
溶接外観検査 & 要 / 不要 & 要 / 不要 &  \\ \hline
非破壊検査 & 要 / 不要 & 要 / 不要 & PT / MT / UT / RT \\ \hline
組立動作確認 & 要 / 不要 & 要 / 不要 &  \\ \hline
\end{tabularx}

\subsection{重点管理寸法・特性}
\begin{tabularx}{\linewidth}{|L{8mm}|L{22mm}|Y|L{22mm}|L{26mm}|L{22mm}|}
\hline
\rowcolor{SpecSoft}\thcell{No.} & \thcell{図面位置} & \thcell{寸法 / 特性} & \thcell{公差} & \thcell{検査方法} & \thcell{備考} \\ \hline
1 &  &  &  &  &  \\ \hline
2 &  &  &  &  &  \\ \hline
3 &  &  &  &  &  \\ \hline
4 &  &  &  &  &  \\ \hline
5 &  &  &  &  &  \\ \hline
\end{tabularx}

\begin{SpecNote}
検査成績書が必要な場合は、測定値、使用測定器、検査日、検査者を記録すること。
\end{SpecNote}

\subsection{検査・判定に関する補足}
\begin{tabularx}{\linewidth}{|L{38mm}|Y|}
\hline
\labelcell{抜取 / 全数} & 全数 / 抜取 / 重要寸法のみ全数 \\ \hline
\labelcell{測定器指定} & 指定なし / 三次元測定機 / 粗さ計 / 硬度計 / その他 \\ \hline
\labelcell{判定保留時の扱い} & 発注者確認後に処置すること。 \\ \hline
\end{tabularx}

\PageBreak

% =========================================================
\section{提出物・納入条件}

\subsection{提出書類・納入物}
\begin{tabularx}{\linewidth}{|Y|L{22mm}|L{34mm}|Y|}
\hline
\rowcolor{SpecSoft}\thcell{提出物} & \thcell{要否} & \thcell{提出タイミング} & \thcell{備考} \\ \hline
見積書 & 要 / 不要 & 発注前 &  \\ \hline
工程表 & 要 / 不要 & 製作前 &  \\ \hline
材料証明書 & 要 / 不要 & 納入時 &  \\ \hline
熱処理証明書 & 要 / 不要 & 納入時 &  \\ \hline
表面処理証明書 & 要 / 不要 & 納入時 &  \\ \hline
検査成績書 & 要 / 不要 & 納入時 &  \\ \hline
非破壊検査記録 & 要 / 不要 & 納入時 &  \\ \hline
完成写真 & 要 / 不要 & 出荷前 / 納入時 &  \\ \hline
\end{tabularx}

\subsection{梱包要求}
\begin{tabularx}{\linewidth}{|Y|Y|}
\hline
\cb\ 輸送中の傷、打痕、錆、変形を防止 & \cb\ 精密面、シール面、摺動面を保護 \\ \hline
\cb\ ねじ部、突起部、薄肉部を保護 & \cb\ 必要に応じ防錆油、防錆紙、乾燥剤を使用 \\ \hline
\cb\ 重量物は安全に荷扱い可能な荷姿 & \cb\ 混載時は部品識別を明確化 \\ \hline
\end{tabularx}

\subsection{識別・納入条件}
\begin{tabularx}{\linewidth}{|L{35mm}|Y|}
\hline
\labelcell{部品識別} & 品名 / 図面番号 / Rev. / 数量を表示 \\ \hline
\labelcell{個別刻印} & 要 / 不要 / 指定部のみ \\ \hline
\labelcell{ロット管理} & 要 / 不要 \\ \hline
\labelcell{納品書記載} & 品名、図面番号、Rev.、数量を記載 \\ \hline
\labelcell{納入場所} &  \\ \hline
\end{tabularx}

\PageBreak

% =========================================================
\section{品質保証・変更管理}

\subsection{不適合・疑義発生時の報告対象}
\begin{tabularx}{\linewidth}{|Y|Y|}
\hline
\cb\ 図面または仕様を満足しない寸法、形状、外観 & \cb\ 材料、処理、溶接、組立に関する不適合 \\ \hline
\cb\ 納期に影響する問題 & \cb\ 図面・仕様の解釈に疑義がある事項 \\ \hline
\cb\ 製作上、仕様変更が必要と考えられる事項 & \cb\ 補修、再製作、特別採用が必要な事項 \\ \hline
\end{tabularx}

\begin{SpecNote}
不適合品の補修、再製作、特別採用は、発注者の承認なく実施しないこと。
\end{SpecNote}

\subsection{事前承認を要する変更}
\begin{tabularx}{\linewidth}{|Y|Y|}
\hline
材料の変更 & 加工方法の大幅な変更 \\ \hline
溶接方法または溶接材料の変更 & 熱処理または表面処理条件の変更 \\ \hline
外注先、再委託先の変更 & 図面・仕様からの逸脱 \\ \hline
納入形態の変更 & 補修方法の変更 \\ \hline
\end{tabularx}

\subsection{見積時確認事項}
\begin{tabularx}{\linewidth}{|Y|L{42mm}|Y|}
\hline
\rowcolor{SpecSoft}\thcell{確認項目} & \thcell{回答} & \thcell{備考} \\ \hline
製作可否 & 可 / 不可 / 条件付き可 &  \\ \hline
希望納期対応 & 可 / 不可 &  \\ \hline
材料手配 & 可 / 不可 &  \\ \hline
表面処理対応 & 自社 / 外注 / 不可 &  \\ \hline
溶接対応 & 自社 / 外注 / 不可 &  \\ \hline
検査成績書提出 & 可 / 不可 &  \\ \hline
\end{tabularx}

\subsection{受入条件}
\begin{enumerate}
  \item 図面および本仕様書の要求を満足していること。
  \item 必要な検査および提出書類が揃っていること。
  \item 外観、寸法、機能に重大な不適合がないこと。
  \item 梱包、識別、数量、Rev. が発注内容と一致していること。
  \item 不明点、逸脱、変更について事前承認が得られていること。
\end{enumerate}

\PageBreak

% =========================================================
\section{特記事項・承認}

\subsection{案件固有の特記事項}
\begin{tabularx}{\linewidth}{|L{42mm}|Y|}
\hline
\labelcell{使用環境} &  \\ \hline
\labelcell{取付相手 / インターフェース} &  \\ \hline
\labelcell{荷重条件 / 使用条件} &  \\ \hline
\labelcell{温度条件 / 耐食性要求} &  \\ \hline
\labelcell{禁止事項} &  \\ \hline
\labelcell{その他} & \rule{0pt}{10mm} \\ \hline
\end{tabularx}

\subsection{最低記入項目チェック}
\begin{tabularx}{\linewidth}{|L{30mm}|Y|L{16mm}|}
\hline
\rowcolor{SpecSoft}\thcell{分類} & \thcell{必須項目} & \thcell{確認} \\ \hline
基本情報 & 品名、図面番号、Rev.、数量、納期 & \cb \\ \hline
技術情報 & 図面、材質、表面処理、熱処理、一般公差 & \cb \\ \hline
製作範囲 & 加工、溶接、組立、材料手配の要否 & \cb \\ \hline
品質要求 & 検査項目、検査成績書の要否、重点管理寸法 & \cb \\ \hline
納入条件 & 梱包、識別、提出書類 & \cb \\ \hline
変更管理 & 代替材、仕様逸脱、補修時の承認要否 & \cb \\ \hline
\end{tabularx}

\subsection{承認欄}
\begin{tabularx}{\linewidth}{|Y|Y|Y|Y|}
\hline
\labelcell{作成} & \labelcell{確認} & \labelcell{承認} & \labelcell{受領確認} \\ \hline
\rule{0pt}{22mm}日付： & 日付： & 日付： & 日付： \\ \hline
\end{tabularx}

\end{document}
